\documentclass[11pt, a4paper]{article}

% --- Encoding & Fonts ---
\usepackage[utf8]{inputenc}
\usepackage[T1]{fontenc}

% --- Layout & Formatting ---
\usepackage[margin=1in]{geometry}
\usepackage{booktabs}   % Professional tables
\usepackage{array}      % >{\raggedright\arraybackslash}p{...} columns
\usepackage{xcolor}     % Colors
\usepackage{graphicx}

% --- Math Packages ---
\usepackage{amsmath, amssymb, amsfonts}
\usepackage{amsthm}
\usepackage[most]{tcolorbox}  % colored callout boxes (boxed `question' env)

% --- Bibliography ---
% Use BibTeX via \bibliography{...} (not biblatex/\printbibliography)
\usepackage{cite}

% --- Hyperlinks (Load last) ---
\usepackage[colorlinks=true, linkcolor=blue, citecolor=blue, urlcolor=blue]{hyperref}
\usepackage[nameinlink,capitalise]{cleveref}

% --- Theorem Environments Setup ---
\theoremstyle{definition}
\newtheorem{thm}{Theorem}[section]
\newtheorem{lem}[thm]{Lemma}
\newtheorem{rem}[thm]{Remark}
\newtheorem{example}[thm]{Example}
\newtheorem{prop}[thm]{Proposition}
\newtheorem{defi}[thm]{Definition}
\newtheorem{question}[thm]{Question}
\newtheorem{hypothesis}[thm]{Hypothesis}

% --- Boxed "question" environment ---
\tcolorboxenvironment{question}{
  enhanced jigsaw,
  breakable,
  colback=blue!5,
  colframe=blue!55,
  boxrule=0.6pt,
  arc=2pt,
  left=8pt, right=8pt, top=6pt, bottom=6pt,
  before skip=10pt, after skip=10pt,
}

% --- cleveref names ---
% amsthm environments share the `thm' counter, so cleveref cannot infer their
% names; declare each one it is asked to typeset.
\crefname{example}{example}{examples}
\Crefname{example}{Example}{Examples}
\crefname{prop}{Proposition}{Propositions}
\Crefname{prop}{Proposition}{Propositions}
\crefname{defi}{Definition}{Definitions}
\Crefname{defi}{Definition}{Definitions}
\crefname{lem}{Lemma}{Lemmas}
\Crefname{lem}{Lemma}{Lemmas}
\crefname{rem}{Remark}{Remarks}
\Crefname{rem}{Remark}{Remarks}

% --- Custom Commands ---
\newcommand{\R}{\mathbb{R}}
\newcommand{\Dist}{\Delta}   % \Dist(S): probability distributions on the set S
\newcommand{\noi}{\noindent}
\newcommand{\ra}{\rightarrow}
% \AA, \BB, ..., \ZZ -> \mathbb{A}, ..., \mathbb{Z}
% NB: overrides kernel \AA (Å) and \SS (ß); neither is used in this project.
\makeatletter
\@for\@bbletter:=A,B,C,D,E,F,G,H,I,J,K,L,M,N,O,P,Q,R,S,T,U,V,W,X,Y,Z\do{%
  \expandafter\edef\csname\@bbletter\@bbletter\endcsname{%
    \noexpand\mathbb{\@bbletter}}%
}
\makeatother
%brackets
\newcommand{\anbr}[1]{\left\langle #1 \right\rangle}
\newcommand{\crbr}[1]{\left \{ #1 \right \} }
\newcommand{\smbr}[1]{\left( #1 \right) }
\newcommand{\sqbr}[1]{\left [ #1 \right ] }
% \c is \mathcal in math mode but keeps its kernel meaning (cedilla accent,
% e.g. Fran\c{c}a in the bibliography) in text mode.
\let\cedillaaccent\c
\renewcommand{\c}{\ifmmode\expandafter\mathcal\else\expandafter\cedillaaccent\fi}

%commenting
\newcommand{\milton}[1]{\textcolor{blue!55}{[milton: #1]}}

% Figures are inline, uncaptioned and unnumbered: they render exactly where
% they are written, including inside an itemize/enumerate, where a float
% cannot go.  Optional argument = width as a fraction of the current line.
\newcommand{\inlinefig}[2][1.0]{%
  \par\vspace{0.8em}%
  \noindent\makebox[\linewidth][c]{\includegraphics[width=#1\linewidth]{#2}}%
  \par\vspace{0.8em}%
}

\title{A Controlled Study of Pure Monte Carlo Tree Search for Symbolic Regression}
\author{After conversations with GPT}
\date{\today}

\begin{document}
\maketitle
\setlength{\emergencystretch}{3em}  % absorb long \texttt{} tokens without margin overflow

\section{Setup}
\label{sec:pure-mcts-study}

\subsection{Problem}

Data \((x_i,y_i)_{i=1}^n\) are generated by a fixed unknown expression \(h\),
\(y_i=h(x_i)\).  Symbolic regression searches a discrete set \(\Pi\) of
expression \emph{structures}; a structure \(\pi\in\Pi\) carries free
coefficients \(\beta\), fitted by least squares,
\[
    \hat\beta_\pi
    =
    \arg\min_\beta \frac1n\sum_{i=1}^n\bigl(y_i-\pi_\beta(x_i)\bigr)^2,
\]
and is scored by the coefficient of determination of that fit, clipped to
\([-1,1]\) and written \(R(\pi)\) \eqref{eq:code-reward}.  The clip is
load-bearing rather than cosmetic: fitting \(\beta x\) to \texttt{lin\_C} of
\cref{tab:family} returns \(R^2\approx-2.9\times10^{12}\), and one such
completion would otherwise decide the mean value of a whole subtree.  The
inner fit is part of the score, so \(R\) is a function of the structure alone
and the symbolic search is \(\pi^*=\arg\max_{\pi\in\Pi}R(\pi)\).  The
motivating instance is a quadratic target \(h(x)=c_0+c_1x+c_2x^2\), with
structures ranging over the subsets of \(\{1,x,x^2\}\).

The structure search becomes \emph{subset selection}.  The purpose of the
study is to characterize the finite-budget limitations of pure Monte Carlo
tree search (MCTS) in grammar-based symbolic regression.  The
neural-guidance question is scientifically separate: before asking whether a
network improves search, one must determine what information the underlying
pure-search procedure estimates and when that information is aligned with
symbolic optimization.

\subsection{Values and failure modes}
\label{subsec:values-failure-modes}

A state \(s\) is a partial derivation, \(\c{A}(s)\) its set of legal
actions, and \(T:\c{S}\times\c{A}\ra\c{S}\) the deterministic transition.
Let
\[
    \c{E}(s)
    =
    \crbr{\pi : \pi \text{ is a terminal expression completable from } s},
\]
so that \(\Pi=\c{E}(s_0)\) is the structure set of \S\ref{sec:pure-mcts-study}
and a terminal expression \(\pi\) is itself a state, with
\(\c{E}(\pi)=\{\pi\}\).  The best-achievable value is
\begin{equation}
    V^*(s)
    =
    \max_{\pi\in\c{E}(s)} R(\pi),
    \label{eq:vstar}
\end{equation}
where \(R(\pi)\) is the terminal symbolic-regression reward.  The uniform
completion policy
\[
    q(a\mid s)=\frac{1}{|\c{A}(s)|},
    \qquad a\in\c{A}(s),
\]
induces, because transitions are deterministic, a distribution
\(q(\cdot\mid s)\in\Dist(\c{E}(s))\) over completions; the mean-rollout
value is
\begin{equation}
    V^q(s)
    =
    \mathbb E_{\pi\sim q(\cdot\mid s)}\bigl[R(\pi)\bigr].
    \label{eq:vroll}
\end{equation}
Every target below is exactly expressible in the grammar, so \(V^*(s_0)=1\),
and the exact-completion mass is
\begin{equation}
    \rho(s)
    =
    \Pr_{\pi\sim q(\cdot\mid s)}\bigl[R(\pi)\ge1-\tau\bigr],
    \qquad
    \tau=10^{-6},
    \label{eq:rare-mass}
\end{equation}
the probability that a uniform completion of \(s\) fits the target exactly.
The threshold is the \emph{global} optimum \(V^*(s_0)=1\), not the local
\(V^*(s)\), so it is the same at every state and \(\rho\) obeys the same
backward induction as \(V^q\).  A single tolerance \(\tau\) serves here and in
the exact-recovery indicator \eqref{eq:exact-recovery-indicator}; on the six
targets that carry the study every exact fit returns \(R=1\) to the last bit,
so no reported value of \(\rho\) depends on it (\cref{rem:conditioning}).

\begin{question}
Which mechanisms make pure MCTS fail at finite budget?
\end{question}

\noi Two candidate mechanisms:
\begin{enumerate}
    \item \emph{Value-semantic failure}: for sibling states \(s_1,s_2\),
    \[
        V^q(s_1)>V^q(s_2)
        \quad\text{while}\quad
        V^*(s_1)<V^*(s_2).
    \]
    Mean-based search then confidently exploits the wrong branch: a shallow
    decoy yields a high expected reward, while the branch containing the
    true optimum is penalized by its abundance of poor completions.
    \item \emph{Rare-event failure}: \(\rho(s)\ll B^{-1}\) at
    budget \(B\).  Search may explore the correct branch yet exhaust the
    budget before sampling a near-optimal leaf.
\end{enumerate}

\noi Therefore \(V^q(s)\), \(V^*(s)\) and \(\rho(s)\) are the
three primary metrics of every experiment below.
\Cref{tab:failure-modes} summarizes the two mechanisms.

\begin{table}[thbp!]
    \centering
    \footnotesize
    \renewcommand{\arraystretch}{1.35}
    \begin{tabular}{@{}
        >{\raggedright\arraybackslash}p{2.2cm}
        >{\raggedright\arraybackslash}p{3.4cm}
        >{\raggedright\arraybackslash}p{4.3cm}
        >{\raggedright\arraybackslash}p{4.5cm}
        @{}}
        \toprule
        \textbf{Failure mechanism} & \textbf{Formal condition} &
        \textbf{Intuitive cause} & \textbf{Implication for search} \\
        \midrule
        Value-semantic failure &
        \(V^q(s_1) > V^q(s_2)\) while \(V^*(s_1) < V^*(s_2)\)
        for sibling states \(s_1, s_2\) &
        A suboptimal branch has the higher mean rollout reward: a shallow
        decoy looks better &
        Search confidently exploits the wrong branch and starves the branch
        containing the true global optimum. \\
        \addlinespace
        Rare-event failure &
        \(\rho(s) \ll B^{-1}\) &
        Near-optimal completions carry a negligible fraction of the rollout
        probability mass. &
        Search explores the correct branch but exhausts the budget \(B\)
        before sampling a near-optimal leaf. \\
        \bottomrule
    \end{tabular}
    \caption{The two finite-budget failure mechanisms of mean-based MCTS.}
    \label{tab:failure-modes}
\end{table}

\begin{rem}
The optimization objective depends on \(V^*\), whereas conventional UCT uses
empirical mean returns as its exploitation statistic \cite{kocsis2006bandit};
asymptotic consistency does not imply effective finite-budget search, and
adverse finite-depth behavior is known for UCT \cite{coquelin2007bandit}.
The goal in symbolic regression is one exceptional expression, not a policy
with high average return: best-focused Monte Carlo methods
\cite{cazenave2009nested,cazenave2013montecarlo}, risk-seeking symbolic
regression \cite{petersen2021deep}, and recent MCTS systems with learned
proposals or extreme-value allocation
\cite{kamienny2023deep,huang2025improving} all respond to this mismatch.
\end{rem}

Experiment~1 measures these quantities exactly on a controlled family.
Experiment~2 tests whether they predict the finite-budget behavior of pure
MCTS.

\subsection{Notation}
\label{subsec:notation}

Three families of symbol are easy to confuse and are kept apart throughout:
math italic \(c_p\) for the coefficients of the \emph{target}, typewriter
\texttt{C0},~\texttt{C1},~\texttt{C2} for the \emph{grammar tokens} whose
values are fitted, and \(C\) for a distinguished \emph{state}.

\begin{center}
\footnotesize
\renewcommand{\arraystretch}{1.2}
\begin{tabular}{@{}lll@{}}
    \toprule
    Symbol & Type & Meaning \\
    \midrule
    \(s\in\c{S}\)               & state    & sentential form in the length-\(L\) buffer \\
    \(s_0\)                     & state    & root, the one-symbol form \((S)\) \\
    \(\pi\in\Pi\)               & state    & terminal expression; \(\Pi=\c{E}(s_0)\) \\
    \(\c{E}(s)\)                & set of states & terminals completable from \(s\) \\
    \(a=\mathrm{pos}\cdot P+\mathrm{prod}\) & action & flattened \((\mathrm{pos},\mathrm{prod})\) index \\
    \(\c{A}(s)\)                & set of actions & the legal mask at \(s\), \eqref{eq:action-mask} \\
    \(T(s,a)\)                  & map \(\c{S}\times\c{A}\ra\c{S}\) & deterministic transition \\
    \(w_0,\ldots,w_{P-1}\)      & symbol strings & the productions, in code order \\
    \(S\), \(A\)                & nonterminals of \(G\) & italic; calligraphic \(\c{S},\c{A}\) are the state and action sets \\
    \(R(\pi)\)                  & \(\R\)   & \(\operatorname{clip}(R^2,-1,1)\) after the inner fit, \eqref{eq:code-reward} \\
    \(V^*,V^q,\rho\) & \(\c{S}\ra\R\) & \eqref{eq:vstar}, \eqref{eq:vroll}, \eqref{eq:rare-mass} \\
    \(C,D_1,D_2,D_3\)           & states   & the four children of \(s_0\), \cref{subsec:difficulty} \\
    \(c_0,c_1,c_2\)             & \(\R\)   & the target's chosen coefficients, \eqref{eq:target} \\
    \texttt{C0},\texttt{C1},\texttt{C2} & grammar tokens & free constants fitted by \texttt{lmfit} \\
    \(B\), \(K\)                & \(\NN\)  & terminal-evaluation budget; independent completions \\
    \(\tau\)                    & \(\R\)   & exactness tolerance \(10^{-6}\), \eqref{eq:rare-mass} \\
    \bottomrule
\end{tabular}
\end{center}

\noi The reachable states form a directed acyclic graph rather than a tree:
an action names a buffer position, so expanding the two nonterminals of
\texttt{+ S S} in either order reaches the same form.  Experiment~1 therefore
runs a backward induction over \(4{,}898\) \emph{states}, whereas the search
tree Experiment~2 builds is a tree over \emph{action sequences}, in which one
form may occupy several nodes.

\section{The derivation MDP in the code}
\label{sec:derivation-mdp}

\begin{defi}[The grammar-derivation MDP]
\label{def:derivation-mdp}
Fix a context-free grammar \(G\) with start symbol \(S\), production set
\(\{w_0,\ldots,w_{P-1}\}\) indexed in code order, and a buffer length \(L\)
(\texttt{max\_len}).  The environment is the deterministic undiscounted MDP:
\begin{itemize}
    \item \emph{State.}  A sentential form \(s\in\Sigma^*\) with
    \(\lvert s\rvert\le L\), stored as a fixed-length array of \(L\) symbols
    together with \(\lvert s\rvert=\)\,\texttt{real\_state\_len}; cells beyond
    \(\lvert s\rvert\) hold a pad symbol.  The observation space is
    \(\texttt{MultiDiscrete}([\,\lvert\Sigma\rvert+1\,]^{L})\), so the pad
    symbol is observable and the agent sees how much buffer
    remains.\footnote{Load-bearing, not cosmetic: the policy/value net is a
    fixed-width MLP over a one-hot encoding of the whole buffer, \(L(\lvert
    \Sigma\rvert+1)\) inputs built by \texttt{SymRegPolicyValueNet.\_encode}.
    \texttt{real\_state\_len} is never itself passed to the network, only
    \texttt{state} is, so a pad symbol in the vocabulary is the only way a
    fixed-shape input still carries \(|s|\): the network recovers it by
    counting non-pad slots.}
    \inlinefig{figures/state_buffer.png}
    \item \emph{Start.}  \(s_0=(S)\), a single nonterminal, \(\lvert s_0\rvert=1\).
    \item \emph{Action.}  An action must say \emph{which} cell of the buffer
    to rewrite and \emph{which} production to write there, so the action set is
    the set of pairs
    \((\mathrm{pos},\mathrm{prod})\in\{0,\ldots,L-1\}\times\{0,\ldots,P-1\}\).
    There are \(LP\) such pairs, and \texttt{Discrete(L*P)} indexes them by one
    integer counted from zero, whence
    \begin{equation}
        a\in\{0,\ldots,LP-1\},
        \qquad
        a=\mathrm{pos}\cdot P+\mathrm{prod},
        \qquad
        (\mathrm{pos},\mathrm{prod})=\operatorname{divmod}(a,P).
        \label{eq:action-decode}
    \end{equation}
    The \(-1\) is zero-indexing of \(LP\) cells and nothing more; at \(L=12\),
    \(P=4\) the largest action is \(47\).
    \inlinefig{figures/action_index.png}
    \item \emph{Legal set.}  \(a\in\c{A}(s)\) iff
    \begin{equation}
        \mathrm{pos}<\lvert s\rvert,
        \quad
        s[\mathrm{pos}]\text{ is a nonterminal},
        \quad
        \mathrm{prod}\in\mathrm{proddict}[s[\mathrm{pos}]],
        \quad
        \lvert s\rvert+\lvert w_{\mathrm{prod}}\rvert-1<L.
        \label{eq:action-mask}
    \end{equation}
    \item \emph{Transition.}  \(T(s,a)\) splices \(w_{\mathrm{prod}}\) over cell
    \(\mathrm{pos}\), shifts the tail right, and re-pads.  It is a function, not
    a kernel.
    \item \emph{Reward.}  \(R\equiv0\) at every nonterminal state.  At a
    terminal, the symbol string is read as a prefix expression, its
    \texttt{C*} tokens are fitted by \texttt{lmfit}, and the reward is the
    clipped coefficient of determination
    \begin{equation}
        R(\pi)=\operatorname{clip}\bigl(R^2,-1,1\bigr),
        \label{eq:code-reward}
    \end{equation}
    with \(-1\) returned for any unparseable expression or failed solve.
    \item \emph{Discount.}  \(\gamma=1\); exactly one reward per episode is
    nonzero, so the return of an episode is its terminal \(R^2\).
\end{itemize}
\end{defi}

\inlinefig{figures/game_mdp.png}

Two features of \eqref{eq:action-decode}--\eqref{eq:action-mask} matter for
everything downstream.  First, the action space is nominally \(LP\) but almost
entirely illegal: on the additive instance at \(L=12\), \(LP=48\) while the
legal set averages \(3.38\) and never exceeds \(16\) over all \(4{,}651\)
nonterminal forms.  The mask, not the action space, is what search branches on.
Second, because an action names a position, the uniform completion policy \(q\)
of \eqref{eq:vroll} is uniform over the flattened
\((\mathrm{pos},\mathrm{prod})\) mask, \emph{not} uniform over the productions
of one designated nonterminal.  The two measures differ as soon as a form holds
more than one nonterminal, and \(V^q\) inherits whichever is used, so which one
is meant must be stated.


\subsection{Termination}
\label{subsec:termination}

The environment has no step limit, no truncation, and no stop action.  An
episode ends when the sentential form runs out of nonterminals, and not before.
Concretely, \texttt{step} has three exits, and only the first can fire under
masked play.

\paragraph{Exit 1: normal termination.}
\(\texttt{terminated}=\neg\,\texttt{\_has\_nonterms}(s)\).  The finished string
is decoded, converted to infix, fitted, and scored by \eqref{eq:code-reward}.
This is how every episode in every experiment ends.

\paragraph{Exit 2: invalid action.}
A defensive branch returns \(R=-1\) with \(\texttt{terminated}=\)~true when the
decoded action is out of range, addresses a terminal cell, or would overflow the
buffer.  Calling \texttt{step} off-mask on \(s_0\) reproduces it exactly.  Under
masked play it is unreachable; it exists so that a masking bug fails loudly
rather than looping.

\paragraph{Exit 3: dead end.}
A form retains a nonterminal but \(\c{A}(s)=\emptyset\).  Then \texttt{step} is
never called at all: \texttt{MCTS.\_rollout\_value} breaks out of its loop and
\emph{discards} the sample rather than scoring it.
\Cref{prop:no-dead-ends} shows this exit cannot fire.

\begin{prop}[Length bound]
\label{prop:length-bound}
Every reachable form satisfies \(\lvert s\rvert\le L-1\).  In particular every
terminal expression has at most \(L-1\) symbols.
\end{prop}

\begin{proof}
Induction on the number of actions.  \(\lvert s_0\rvert=1\le L-1\).  If \(s'\)
follows \(s\) by a legal action with production \(w\), the mask condition
\(\lvert s\rvert+\lvert w\rvert-1<L\) in \eqref{eq:action-mask} states exactly
that \(\lvert s'\rvert<L\), hence \(\lvert s'\rvert\le L-1\).
\end{proof}

The bound is attained, and the strict inequality in \eqref{eq:action-mask} is
load-bearing: at \(L=12\) the longest terminals have \(11\) symbols, and the
exact expression \texttt{+ C0 + * C1 x * C2 * x x} is admissible so that
\(V^*(s_0)=1\).  At \(L=11\) it is excluded, and \(V^*(s_0)\) then depends on
the target: it stays at \(1\) for every linear target, which needs only five
symbols, but falls to \(0.997213\) on \texttt{quad\_B} and \(0.527856\) on
\texttt{quad\_A} (\cref{sec:family}).  A buffer must be set one longer than the
longest expression one intends to be reachable.

\begin{prop}[No dead ends]
\label{prop:no-dead-ends}
Suppose every nonterminal \(A\) of \(G\) admits a production \(A\to w\) with
\(\lvert w\rvert=1\), an \emph{escape production}.  Then \(\c{A}(s)\neq\emptyset\)
for every reachable \(s\) containing a nonterminal, so Exit 3 never fires.
\end{prop}

\begin{proof}
Let \(s[i]=A\) be a nonterminal and \(A\to w\) an escape production.  The first
three conditions of \eqref{eq:action-mask} hold at that pair by construction.
For the fourth, \(\lvert s\rvert+\lvert w\rvert-1=\lvert s\rvert\le L-1<L\) by
\cref{prop:length-bound}.  The action is legal, so the mask is nonempty.
\end{proof}

Both shipped grammars satisfy the hypothesis through \(S\to\)~\texttt{C0}; an
exhaustive census finds \(0\) dead ends among the \(4{,}898\) reachable forms.
The hypothesis is a genuine constraint rather than a formality.  Removing the
constant production, or capping \(L\) below the shortest completion of some
reachable form, reintroduces dead ends --- and since the rollout discards them
instead of scoring them, they would bias \(V^q\) silently rather than raise an
error.

\inlinefig{figures/deadend_causes.png}

\begin{rem}[Parity]
Each action changes the length by \(\lvert w\rvert-1\in\{0,2,4\}\), all even,
and \(\lvert s_0\rvert=1\).  Every reachable form therefore has odd length, and
terminals can only hold \(1,3,5,7,9,11\) symbols.  The census agrees exactly:
\(1,2,5,14,48,177\) terminals at those six lengths.
\end{rem}

\inlinefig{figures/game_termination.png}

\section{A minimal controlled family}
\label{sec:family}

The MDP of \cref{def:derivation-mdp} is now held fixed --- grammar
\texttt{ADDITIVE\_GRAMMAR}, buffer \(L=12\) --- and the target the terminal
expression is fitted against is the only thing that varies.

\begin{defi}[Target family]
\label{def:target-family}
A \emph{target} is a polynomial together with the grid it is sampled on,
\begin{equation}
    y(x)=\sum_{p=0}^{d}c_p\,x^p,
    \qquad
    x_1,\ldots,x_n\text{ equally spaced in }[-1,1],
    \quad n=41,
    \label{eq:target}
\end{equation}
with \(d\le2\) and coefficients \emph{chosen}.  The family has eight
members: four linear, separated by where the zero \(-c_0/c_1\) falls relative
to the domain, and four quadratic, separated by the vertex \(x_v=-c_1/2c_2\)
and the discriminant \(\operatorname{disc}=c_1^2-4c_0c_2\).\footnote{\(x_v\)
says where the curve turns and \(\operatorname{disc}\) whether it meets the
axis at all (\(\operatorname{disc}>0\) two real zeros, \(<0\) none), so the
pair fixes the shape on \([-1,1]\): of the four members only \texttt{quad\_B}
is monotone there, and only \texttt{quad\_D} crosses zero inside the domain.
The word \emph{root} is reserved throughout for the MDP root \(s_0\).}
\end{defi}

\inlinefig{figures/family_targets.png}


\subsection{Where the difficulty is}
\label{subsec:difficulty}


Exactly \(\lvert\c{A}(s_0)\rvert=4\) actions are legal at the root, one per
production, and they have four distinct successors.  We name those four
successor \emph{states} after the production that creates them,
\begin{equation}
    D_1=\texttt{C0},
    \qquad
    D_2=\texttt{* C1 x},
    \qquad
    D_3=\texttt{* C2 * x x},
    \qquad
    C=\texttt{+ S S},
    \label{eq:root-decoys}
\end{equation}
the same four states already previewed in \cref{subsec:notation}, so that
\(D_i,C\in\c{S}\) and every value written against them --- \(V^*(C)\),
\(V^q(C)\), \(\rho(C)\), \(R(D_i)\) --- is the state function of
\S\ref{subsec:values-failure-modes} evaluated at that state.  Because the four
root actions are in bijection with these four children, we also write
\(N(s_0,C)\) for the visit count of the edge \(s_0\ra C\); no separate symbol
for a root action is needed.

\paragraph{Three of the four are already decided.}
Three of the four productions carry no nonterminal on the right, so \(D_1\),
\(D_2\) and \(D_3\) are already terminal: each is a one-symbol derivation with
\(V^*(D_i)=V^q(D_i)=R(D_i)\), one deterministic number.  The fourth child
\(C\) is the only one that continues, and the only one whose subtree contains
an exact expression: \(V^*(C)=1\) for every target.  Root value-semantic
failure is therefore a single inequality,
\begin{equation}
    \max_i R(D_i)>V^q(C)
    \qquad\text{while}\qquad
    \max_i R(D_i)<V^*(C)=1,
    \label{eq:inversion}
\end{equation}
and nothing has to be added to the environment to produce it.

\inlinefig{figures/family_difficulty.png}

Both sides of \eqref{eq:inversion} are tabulated in \cref{tab:family}.  Only
\texttt{lin\_A} and \texttt{quad\_D} satisfy it.  For \texttt{lin\_D} and
\texttt{quad\_C} the margin \(\max_iR(D_i)-V^q(C)\) is also positive, but there
the one-action terminal is itself exact --- for \texttt{quad\_C} only to
within round-off, \cref{rem:conditioning} --- so preferring it is correct and
no trap exists.  The second inequality of \eqref{eq:inversion}, not the margin
alone, separates the two cases.

\begin{table}[thbp!]
    \centering
    \footnotesize
    \renewcommand{\arraystretch}{1.15}
    \begin{tabular}{@{}llrrrrrrl@{}}
        \toprule
        Target & \(y(x)\)
        & \(R(D_1)\) & \(R(D_2)\) & \(R(D_3)\)
        & \(V^q(C)\) & \(\rho(C)\) & margin & \\
        \midrule
        \texttt{lin\_A}  & \(0.5+2x\)              & \(0\) & \(0.8214\)  & \(-0.0793\) & \(0.5646\) & \(0.3945\) & \(+0.2569\) & trap \\
        \texttt{lin\_B}  & \(5+x\)                 & \(0\) & \(-1\)      & \(-1\)      & \(0.1341\) & \(0.3945\) & \(-0.1341\) & --- \\
        \texttt{lin\_C}  & \(1000+0.001x\)         & \(0\) & \(-1\)      & \(-1\)      & \(0.1341\) & \(0.3945\) & \(-0.1341\) & --- \\
        \texttt{lin\_D}  & \(2x\)                  & \(0\) & \(1\)       & \(0\)       & \(0.5924\) & \(0.5924\) & \(+0.4076\) & solved at once \\
        \texttt{quad\_A} & \(1-x+2x^2\)            & \(0\) & \(-1\)      & \(-0.0711\) & \(0.3540\) & \(0.1128\) & \(-0.3540\) & --- \\
        \texttt{quad\_B} & \(6-5x+0.5x^2\)         & \(0\) & \(-1\)      & \(-1\)      & \(0.1563\) & \(0.1128\) & \(-0.1563\) & plateau \\
        \texttt{quad\_C} & \(0.001+0.001x+1000x^2\)& \(0\) & \(-1\)      & \(1\)       & \(0.4840\) & \(0.5569\) & \(+0.5160\) & solved at once \\
        \texttt{quad\_D} & \(-0.48+0.4x+2x^2\)     & \(0\) & \(0.0170\)  & \(0.6461\)  & \(0.5106\) & \(0.1128\) & \(+0.1355\) & trap \\
        \bottomrule
    \end{tabular}
    \caption{Exact root quantities for the eight targets under one fixed MDP;
    \(V^*(C)=1\) in every row.  ``Trap'' marks the rows satisfying
    \eqref{eq:inversion}.  The \(\rho(C)\) column is at \(\tau=10^{-6}\); it is
    tolerance-free except on \texttt{lin\_C} and \texttt{quad\_C}
    (\cref{rem:conditioning}).}
    \label{tab:family}
\end{table}

\paragraph{The last term is the expensive one.}
Panel~(a) plots the reward ceiling: the best \(R^2\) any terminal of \(\ell\)
symbols attains.  For \texttt{quad\_B} the ceiling is \(0.997213\) already at
\(\ell=5\), stays there at \(\ell=7\) and \(\ell=9\), and reaches \(1\) only at
\(\ell=11\).  Three of the five actions of an optimal derivation buy
\(0.0028\) of \(R^2\) between them.  The trap of \eqref{eq:inversion} reappears
one level down as a plateau inside the continue branch, not a decoy action at
the root: a fitted \emph{line} through a parabola whose vertex sits at
\(x_v=5\), far outside the domain, already explains \(99.7\%\) of the
variance.

\paragraph{Rarity is a property of the target's support.}
The tree is the same for every target, so the only thing that can vary is
which of its terminals fit.

\begin{lem}[Rarity depends only on the support]
\label{lem:support}
Let \(P=\{p:c_p\neq0\}\) be the support of the target \eqref{eq:target}.  A
terminal \(\pi\) fits it exactly iff \(\pi\) carries an atom for every
\(p\in P\).  Hence \(\rho\) depends on the target only through \(P\), and
\[
    \rho(C)=0.5924,\quad 0.3945,\quad 0.1128
    \qquad\text{for}\qquad
    P=\{1\},\quad \{0,1\},\quad \{0,1,2\}.
\]
\end{lem}

\begin{proof}
The three productions carrying no nonterminal evaluate to \(1\), \(x\) and
\(x^2\), linearly independent on the \(41\)-point grid, and a terminal is a
sum of such atoms.  Least squares over the atoms present is exact iff the
target lies in their span, which for a polynomial of support \(P\) holds iff
every \(p\in P\) is present, and is otherwise strictly inexact.  The predicate
\(R(\pi)=1\) therefore depends on \(\pi\) only through the set of powers it
carries, and \(\rho\), a uniform average of that predicate over completions
\eqref{eq:backward-induction}, inherits the dependence.
\end{proof}

\noi Of the \(247\) terminals, \(146\) carry both \(1\) and \(x\) while only
\(12\) carry all three, so the family offers exactly two rarity levels.
Since completions are
independent, \(K\) of them miss with
probability \((1-\rho)^K\), so \(\Pr(\text{at least one exact})\ge1-\delta\)
exactly when
\begin{equation}
    K\ \ge\ K_{1-\delta}
    =
    \left\lceil\frac{\log\delta}{\log(1-\rho)}\right\rceil,
    \label{eq:hit-budget}
\end{equation}
which is \(6\) completions at \(\rho=0.39\) and \(26\) at \(\rho=0.11\),
for \(\delta=0.05\).  Here \(K\) counts completions drawn from one fixed state,
whereas the budget \(B\) of \eqref{eq:budget-grid} counts terminal evaluations
anywhere in the tree.  \(K\) is what \emph{uniform} sampling from \(C\)
needs; a tree search is bound by it in neither direction, and
\cref{sec:experiment-two} measures the difference.

\paragraph{The four cells.}
Two binary contrasts are therefore available without writing a line of new
environment code: whether \eqref{eq:inversion} holds, and whether exact
completions are common (\(\rho(C)=0.39\)) or rare (\(0.11\)).

\begin{center}
\footnotesize
\begin{tabular}{@{}lll@{}}
    \toprule
     & common, \(\rho(C)=0.39\) & rare, \(\rho(C)=0.11\) \\
    \midrule
    inversion    & \texttt{lin\_A} & \texttt{quad\_D} \\
    no inversion & \texttt{lin\_B} & \texttt{quad\_A} \\
    \bottomrule
\end{tabular}
\end{center}

\noi \texttt{quad\_B} is kept as a fifth condition: no inversion at the root,
but the \(0.997213\) plateau three actions deep.  These five targets carry the
study.  \texttt{lin\_D} is a solved-at-once control, and the remaining two are
excluded for the reason that follows.

\begin{rem}[The two ill-conditioned members]
\label{rem:conditioning}
\texttt{lin\_C} and \texttt{quad\_C} are the only rows of \cref{tab:family}
that disagree with \cref{lem:support}, and both disagree at the level of
double-precision round-off rather than of mathematics.  \texttt{lin\_C}
\(=10^{-3}\cdot\texttt{lin\_B}+999.995\) is an affine image of
\texttt{lin\_B}; \(R^2\) is invariant under affine maps of \(y\) for every
atom set containing \texttt{C0}, and every set without one clips to \(-1\), so
the two are the same problem and their \(V^q(C)\) agree to
\(3\times10^{-11}\).  Its exact fits nevertheless return
\(R=1-5.6\times10^{-12}\), so \(\rho(C)\) collapses to \(0.2817\) as soon as
\(\tau<5.6\times10^{-12}\).  \texttt{quad\_C} fails the other way: the single
atom \texttt{* C2 * x x} returns \(R=1-8.1\times10^{-12}\), which
\(\tau=10^{-6}\) reads as exact --- whence \(R(D_3)=1\) and \(\rho(C)=0.5569\)
in \cref{tab:family} --- and which any \(\tau<8.1\times10^{-12}\) reads as
inexact, restoring \(\rho(C)=0.1128\).  The window admitting both rows is
\([5.6,8.1)\times10^{-12}\), narrower than a factor of \(1.5\), so no tolerance
makes both mean what they say.  On the other six targets every exact fit
returns \(R^2=1\) to the last bit, and every \(\rho\) reported here is
unchanged for all \(\tau\in[0,10^{-6}]\).
\end{rem}

\section{Experiment 1: exact value-semantic audit}
\label{sec:experiment-one}

\begin{question}
Does the mean random-completion return \(V^q\) preserve the ordering induced
by the best-achievable reward \(V^*\), and how rare
(\(\rho\)) is the exact expression within each partial
derivation?
\end{question}

\paragraph{Hypothesis 1.}
The rollout mean ranks the continue branch above every one-action terminal on
\texttt{lin\_B} and \texttt{quad\_A}, and below one of them on \texttt{lin\_A}
and \texttt{quad\_D}:
\begin{align}
    V^q(C) &> \max_i R(D_i)
    && \text{on \texttt{lin\_B} and \texttt{quad\_A}},
    \nonumber\\
    V^q(C) &< \max_i R(D_i)
    && \text{on \texttt{lin\_A} and \texttt{quad\_D}},
    \label{eq:hypothesis-one-order}
\end{align}
with \(V^*(C)=1>\max_iR(D_i)\) throughout, so the second inequality is a
genuine ordering failure.  Independently, the completion budget required for
exact recovery scales as
\begin{equation}
    K_{1-\delta}
    =
    \Theta\left(
        \frac{\log(1/\delta)}{\rho}
    \right)
    \quad\text{as }\rho\rightarrow0.
    \label{eq:hypothesis-one-rarity}
\end{equation}

\paragraph{Procedure.}
No stochastic tree search is used in this experiment.  Enumerate the
mask-reachable forms once --- \(4{,}898\) of them, \(247\) terminal --- and for
each of the four cells compute exactly, at every state \(s\),
\begin{equation}
    V^*(s),\qquad
    V^q(s),\qquad
    \rho(s).
    \label{eq:experiment-one-quantities}
\end{equation}
For a nonterminal state \(s\), define the action selected by the mean-rollout
oracle as
\begin{equation}
    a_q(s)
    =
    \arg\max_a V^q(T(s,a)),
    \label{eq:mean-oracle-action}
\end{equation}
and define its oracle sibling regret by
\begin{equation}
    g_q(s)
    =
    V^*(s)-V^*(T(s,a_q(s))).
    \label{eq:sibling-regret}
\end{equation}
The reachable set is a directed acyclic graph --- every action either lengthens
the form or removes a nonterminal --- so all three quantities follow by one
backward induction over it,
\begin{equation}
    V^*(s)=\max_{a}V^*(T(s,a)),
    \qquad
    V^q(s)=\frac{1}{\lvert\c{A}(s)\rvert}\sum_{a\in\c{A}(s)}V^q(T(s,a)),
    \label{eq:backward-induction}
\end{equation}
and likewise for \(\rho\), with terminal states seeded by their reward.  The
fitter is called once per terminal and cached, so a cell costs \(247\) fits,
not \(4{,}898\).

Two checks guard the census.  The reward of the exact derivation
\texttt{+ C0 + * C1 x * C2 * x x} must come back as \(1\) to fit tolerance in
every cell, and the induced \(V^q(s_0)\) must agree with the mean of \(10^5\)
uniform masked rollouts to within Monte Carlo error.  Failure of either is an
implementation error, not a finding.

\paragraph{Primary outputs.}
Report the exact root quantities of \cref{tab:family} ---
\(R(D_i)\), \(V^*(C)\), \(V^q(C)\), \(\rho(C)\), \(K_{0.95}\) --- and plot:
\begin{enumerate}
    \item the terminal reward distribution over the \(247\) terminals, per cell;
    \item \(V^q(s)\) against \(V^*(s)\), stratified by \(\lvert s\rvert\);
    \item sibling regret \(g_q(s)\) against \(\lvert s\rvert\);
    \item the exact hit curve \(1-(1-\rho)^K\).
\end{enumerate}

Because Experiment~1 is an exact census, it requires no random seeds,
confidence intervals, or significance tests.

\section{Experiment 2: finite-budget pure-MCTS limits}
\label{sec:experiment-two}

\paragraph{Question.}
At matched terminal-evaluation budgets, do the value-order mismatch and
exact-completion rarity identified in Experiment~1 predict the behavior of
actual pure search?

\subsection{Search conditions}

We compare three conditions.

\paragraph{Uniform random search.}
Sample complete derivations independently from the uniform grammar policy and
return the best terminal expression observed.

\paragraph{Mean-UCT.}
At node \(s\), select
\begin{equation}
    a_t
    =
    \arg\max_a
    \left[
        \overline Q_t(s,a)
        +
        \sqrt{
            \frac{2\log N_t(s)}
                 {N_t(s,a)}
        }
    \right],
    \label{eq:mean-uct}
\end{equation}
where \(N_t(s)\) and \(N_t(s,a)\) count the simulations through \(s\) and
through \((s,a)\) after \(t\) simulations, and \(\overline Q_t(s,a)\) is the
empirical mean of terminal returns from simulations passing through
\((s,a)\).  Unvisited actions receive infinite score.  Rewards lie in
\([-1,1]\), so \eqref{eq:mean-uct} applies the UCB1 bonus for a unit-range
reward to a range-\(2\) one; the constant is held fixed across all conditions
rather than tuned, so it cannot by itself explain a difference between them.
Each simulation selects through the current tree, expands one
unvisited edge, performs one uniform random completion, and backs up the
terminal reward by an arithmetic mean.

\paragraph{Max-UCT diagnostic.}
Define
\begin{equation}
    Q_t^{\max}(s,a)
    =
    \max
    \left\{
        R(\pi_b):
        \pi_b\text{ was sampled through }(s,a)
    \right\}.
    \label{eq:qmax}
\end{equation}
Select using
\begin{equation}
    a_t
    =
    \arg\max_a
    \left[
        Q_t^{\max}(s,a)
        +
        \sqrt{
            \frac{2\log N_t(s)}
                 {N_t(s,a)}
        }
    \right].
    \label{eq:max-uct}
\end{equation}
The tree structure, expansion policy, rollout policy, and exploration term
are otherwise identical to Mean-UCT.  This condition is a causal diagnostic
of the mean-versus-maximum distinction, not a claim that naive maximum backup
is universally optimal.  A calibrated extreme-bandit allocation rule is a
natural subsequent comparison \cite{huang2025improving}.

\subsection{Evaluation protocol}

Each search run uses one persistent tree rooted at \(s_0\).  The algorithm
does not commit to a grammar action before the budget is exhausted.  No
network, learned prior, Dirichlet noise, reward clipping, complexity penalty,
or nonlinear optimizer is used.  One simulation contains exactly one
terminal reward evaluation.

Use the terminal-evaluation budgets
\begin{equation}
    \mathcal B
    =
    \{
        16,32,64,128,256,512,
        1024,2048,4096,8192
    \}.
    \label{eq:budget-grid}
\end{equation}
For each of the four benchmark cells and each search condition, use \(100\)
independent seeds.  Run each seed once to \(B_{\max}=8192\) and record all
intermediate checkpoints.  The same base-seed list is used across algorithms.

\subsection{Outcomes}

Let \(\pi_1,\ldots,\pi_B\) be the terminal expressions observed by budget
\(B\).  Define the best-observed reward
\begin{equation}
    R_B^{\max}
    =
    \max_{1\leq b\leq B}R(\pi_b),
    \label{eq:best-observed-reward}
\end{equation}
the simple regret
\begin{equation}
    r_B
    =
    1-R_B^{\max},
    \label{eq:simple-regret}
\end{equation}
and the exact-recovery indicator
\begin{equation}
    Z_B
    =
    \mathbf 1
    \left\{
        R_B^{\max}\geq1-\tau
    \right\}.
    \label{eq:exact-recovery-indicator}
\end{equation}
The budget at which exact recovery first occurs is
\begin{equation}
    B_{\mathrm{exact}}
    =
    \inf\{B:Z_B=1\},
    \label{eq:time-to-recovery}
\end{equation}
written with \(Z\) and \(B\) rather than \(S\) and \(T\), which
\cref{subsec:notation} reserves for the start symbol and the transition map.

Exact recovery is the primary outcome.  Simple regret is secondary because a
structurally wrong expression already achieves small regret: on
\texttt{quad\_B} the plateau of \cref{subsec:difficulty} sits at
\[
    1-0.997213=0.002787,
\]
and on \texttt{lin\_A} the decoy \(D_2=\texttt{* C1 x}\) sits at
\(1-0.821429=0.178571\) without containing a constant term at all.
Thus predictive reward alone cannot distinguish near-perfect approximation
from structural recovery.  This distinction is consistent with symbolic
regression benchmarks that report both predictive and exact-recovery
performance \cite{lacava2021contemporary}.

For mechanism diagnosis, record the continue-visit fraction
\begin{equation}
    F_B
    =
    \frac{N_B(s_0,C)}{N_B(s_0)},
    \label{eq:continue-fraction}
\end{equation}
where \(N_B(s_0)\) and \(N_B(s_0,C)\) are the root and continue-edge visit
counts at budget \(B\), together with the final root recommendation, number
of distinct terminal
sequences, number of distinct canonical atom sets, and the best reward found
through each root action.  The final visit-count recommendation is diagnostic
only; the primary search output is the best terminal expression encountered.

\subsection{Hypotheses}

\paragraph{Hypothesis 2a: value-semantic effect.}
Mean-UCT will allocate less search to the continue branch on the inversion
cells (\texttt{lin\_A}, \texttt{quad\_D}) than on their non-inversion partners
(\texttt{lin\_B}, \texttt{quad\_A}), because its exploitation statistic is
anchored near \(V^q(C)\) while subtree sampling remains close to uniform, and
\(V^q(C)\) lies below \(\max_iR(D_i)\) exactly on those cells.  Max-UCT will be
less sensitive to the inversion after observing high-reward continuations.

\paragraph{Hypothesis 2b: rarity effect.}
For every algorithm, exact recovery will be slower on the quadratic cells
(\(\rho(C)=0.11\)) than on the linear ones (\(0.39\)).  At moderate budgets,
recovery will be organized approximately by the effective rare-event budget
\begin{equation}
    B\rho.
    \label{eq:effective-rare-budget}
\end{equation}

\paragraph{Hypothesis 2c: interaction.}
The largest recovery difference between Max-UCT and Mean-UCT will occur on
\texttt{quad\_D}, the inversion-and-rare cell, where the continue branch has
both a lower rollout mean than a one-action terminal and the small exact mass
\(\rho(C)=0.1128\).

One pre-specified interaction contrast is
\begin{align}
    \Delta_{\mathrm{semantic}}(B)
    &=
    \left[
        P_B^{\mathrm{Max}}-P_B^{\mathrm{Mean}}
    \right]_{\text{inversion}}
    \nonumber\\
    &\quad -
    \left[
        P_B^{\mathrm{Max}}-P_B^{\mathrm{Mean}}
    \right]_{\text{no inversion}},
    \label{eq:semantic-interaction}
\end{align}
where \(P_B^{\mathrm{Max}}\) and \(P_B^{\mathrm{Mean}}\) are the
exact-recovery probabilities \(\Pr(Z_B=1)\) of Max-UCT and Mean-UCT, each
bracket averaged over the two cells at that level of the design.  The
prediction is
\[
    \Delta_{\mathrm{semantic}}(B)>0,
\]
with a larger effect expected on the quadratic (rare) cells.

\subsection{Statistical reporting and interpretation}

For each budget, cell, and algorithm, report the exact-recovery proportion
with a \(95\%\) Wilson interval, median simple regret with a seed-bootstrap
interval, the empirical survival curve
\[
    \Pr(B_{\mathrm{exact}}>B),
\]
and the continue-visit fraction \(F_B\).

The interpretation is mechanism-based:
\begin{itemize}
    \item If Mean-UCT fails while Max-UCT succeeds, the mean-return statistic
    is causally implicated.
    \item If both methods deteriorate primarily on the quadratic cells,
    completion rarity is the dominant limitation.
    \item If random search outperforms Mean-UCT on an inversion cell, the
    tree allocation rule actively suppresses the optimal branch.
    \item If Mean-UCT succeeds despite the exact value-order inversion, UCT
    exploration is sufficient at the tested budgets and the mismatch alone is
    not a complete explanation.
    \item If Max-UCT performs worse on the non-inversion cells, maximum backup is
    overoptimistic and should not be treated as a universal replacement for
    mean backup.
    \item If simple regret is small while exact recovery remains poor, the
    experiment demonstrates a separation between predictive approximation
    and symbolic identification.
\end{itemize}

\subsection{Secondary validation and deferred questions}

After the two primary experiments, a nonconfirmatory validation may repeat the
key comparison on the full seven-production grammar, whose sine and division
productions are pure distractors for these targets: at \(L=12\) it has
\(795{,}437\) reachable forms and \(19{,}438\) terminals against the \(4{,}898\)
and \(247\) used above, so rarity there is a grammar effect rather than a
target effect.  Nonlinear constant fitting, uniform-prior PUCT, learned policy
and value heads, and SRBench evaluation enter only after the controlled
mechanisms are understood.  In particular,
whether a learned policy or value function corrects the pure-search failures
identified here --- or instead learns and amplifies the biased targets
produced by mean-return search --- is deferred, so that the present study
remains causally interpretable.

% =====================================================================
\bibliographystyle{plain}
\bibliography{references}

\end{document}
